\section{Introduction}
Post-hoc feature attributions for tabular models are routinely computed with Kernel SHAP or permutation SHAP \citep{lundberg2017unified} and then read as \emph{the} explanation of a prediction: the top few attributes are quoted in an audit, a credit decision or a model review. Both estimators are randomised and finite-budget, so two calls with different seeds return different vectors. The estimator literature characterises how the error of such estimates falls with the budget \citep{covert2020improving,mitchell2021sampling,musco2024leverage}; what a practitioner has to decide is different: whether \emph{the one run} they have made, at a budget they can afford, supports the statement they are about to make.

We study this decision on tabular data. The statement we evaluate is the one a reader acts on, the \emph{top-$k$ set} (and sign) of the attributes, not the vector error. We measure it against exact Shapley values where the game can be enumerated and against a validated high-budget reference otherwise, and we put it on the same scale as variation that practitioners already live with: retrained models, redrawn backgrounds and chance. Our contributions are:
\begin{itemize}
\item a measurement protocol for the reliability of a single run (population-level adequacy criterion with a stated Clopper--Pearson limitation, instance-cluster bootstrap, three comparison standards), applied to a grid of \numnCells{} (dataset, model) cells (Sections~\ref{sec:method}--\ref{sec:exp}). What is new relative to the estimator literature is the decision-level target (top-$k$ sets at the budgets a user actually chooses, including the library default) and its comparison with other sources of variation, not the error--budget relation, which we only confirm;
\item the conditions under which one run suffices on these cells: budget, dimension (the library default budget is exact for $d\le11$ and a sampling estimate beyond), model, and the margin between the last attribute shown and the first one omitted;
\item two subset checks on library choices (the default \texttt{l1\_reg} on four cells, the explained output on \numnLogitCells{}) that show these choices change what is reproduced, reported as remarks;
\item negative and null results: no gain from averaging repeats over one larger run, no confirmed advantage of a per-instance diagnostic over what exists, and the experiments we did not run (Section~\ref{sec:limits}).
\end{itemize}
Everything in the paper is regenerated from run code (reproduction notes in the repository).
